
\section{有限型拟凝聚、有限表现与凝聚层}

下面是对一个 $A$-\trModule{} $M$ 的三条自然的有限性条件。若 $A$ 是\trNoetherian{}环——你们中几乎所有人终将会关心的也正是这种情形——这三条便是同一回事。

第一条最简单：一个\trModule{}可以是\textbf{\trFinitelyGenerated}的。换言之，存在\trSurjectiveMap{} $A^{\oplus p}\to M\to 0$。

第二条也合理：它可以是\trFinitelyPresented{}的——即它拥有有限多个生成元与有限多个关系。翻译过来就是：存在一个\textbf{有限表现}（finite presentation），亦即一个\trExactSequence
\[
A^{\oplus q}\longrightarrow A^{\oplus p}\longrightarrow M\longrightarrow 0.
\]

\textbf{6.4.1.}\ 第三个概念说实话有点出人意料。我们称一个 $A$-\trModule{} $M$ 是\textbf{\trCoh}（coherent）的，若 (i)~它是\trFinitelyGenerated{}的，且 (ii)~对任何映射 $A^{\oplus p}\to M$（不必是\trSurjectiveMap！），其\trKernel{}都是\trFinitelyGenerated{}的。

\textbf{6.4.2.\ \trhProposition{}。}\ ——\ \emph{若 $A$ 是\trNoetherian{}的，则这三个定义是同一个。}

\emph{证明。}\ 显然，\trCoh{}蕴涵\trFinitelyPresented，而后者又蕴涵\trFinitelyGenerated。反之，设 $M$ 是\trFinitelyGenerated{}的。任取 $\alpha\colon A^{\oplus p}\to M$。则 $\ker\alpha$ 是 $A$ 上某个\trFinitelyGenerated{}\trModule{}的\trSubmodule，从而由\cref{ex:finitely-generated-noetherian-module}它也\trFinitelyGenerated。于是 $M$ 是\trCoh{}的。\hfill$\square$

于是，大多数人尽可以把这三个概念当成同一回事。

\textbf{6.4.3.\ \trhProposition{}。}\ ——\ \emph{诸\trCoh{}的 $A$-\trModule{}构成 $A$-\trModule{}的\trCategory{}的一个\trAbelian\trSubcategory。}

一般情形的证明是散布在 \S6.7 中的几道小练习。除非你格外好奇，否则不必去做。

\emph{在环 $A$ 为\trNoetherian{}时的证明。}\ 回忆我们在 \S6.3 开头的讨论：必须检验三件事：
\begin{itemize}[leftmargin=2.6em,itemsep=2pt,topsep=3pt]
\item[(i)] 零模是\trCoh{}的。
\item[(ii)] \trCoh{}\trModule{}的\trCategory{}对有限和封闭。
\item[(iii)] \trCoh{}\trModule{}的\trCategory{}对\trKernel{}与\trCokernel{}封闭。
\end{itemize}
前两条是显然的。至于 (iii)，设 $f\colon M\to N$ 是\trFinitelyGenerated{}\trModule{}之间的一个映射。则 $\coker f$ 是\trFinitelyGenerated{}的（它是 $N$ 的\trImage），而 $\ker f$ 亦然（它是\trNoetherian{}环上某个\trFinitelyGenerated{}\trModule{}的\trSubmodule，见\cref{ex:finitely-generated-noetherian-module}）。\hfill$\square$

\textbf{6.4.4.}\ \emph{重要注记：\trPID{}上的\trFinitelyGenerated{}\trModule。}\quad 我们在此记下一笔，以备日后查用：\trPID{}上的\trFinitelyGenerated{}\trModule{}都是\trCyclicModule{}（cyclic module）的有限\trDirectSum（例如见 [DF, \S12.1, Thm.~5]）。我们在此提到它，只因为它总得\emph{在某处}被提上一句。

\textbf{6.4.5.}\quad 我们现在把这三类\trModule{}的定义推广到\trQCSheaf{}上。

\textbf{6.4.A.（练习：有限生成满足\trAffineCommunicationLemma{}~5.3.2 的假设）}\quad

(a)~设 $f\in A$，而 $M$ 是一个\trFinitelyGenerated{}的 $A$-\trModule。证明：$M_f$ 是一个\trFinitelyGenerated{}的 $A_f$-\trModule。

(b)~设 $(f_1,\dots,f_n)=A$，且对一切 $i$，$M_{f_i}$ 都是\trFinitelyGenerated{}的 $A_{f_i}$-\trModule。证明：$M$ 是一个\trFinitelyGenerated{}的 $A$-\trModule。

至于有限表现的情形，下面这个结果颇为趁手，而且它的种种版本也适用于其他场合（「代数」版本见引理~8.3.R）。有很长一段时间，我甚至不曾料到这样的事情竟会是真的。

\textbf{6.4.6.}\ \emph{引理：有限表现蕴涵总是有限表现（模版本）。}\ ——\ \emph{设 $M$ 是一个\trFinitelyPresented{}的 $A$-\trModule{}（对从某个有限自由\trModule{}出发的\textbf{某个}\trSurjectiveMap，其\trKernel{}是\trFinitelyGenerated{}的），又设 $\phi\colon A^{\oplus p}\to M$ 是\textbf{任何}一个\trSurjectiveMap。则 $\ker\phi$ 是\trFinitelyGenerated{}的。}

\emph{证明。}\ 选取 $M$ 的一个有限表现：
\[
\bigoplus_{k=1}^{r} Af_k\overset{\beta}{\longrightarrow}\bigoplus_{i=1}^{n} Ax_i\overset{\alpha}{\longrightarrow}M\longrightarrow 0,
\]
于是 $\alpha(x_1),\dots,\alpha(x_n)$ 生成 $M$。把假设中的那个\trSurjectiveMap{}写作 $\phi\colon\bigoplus_{j=1}^{p}Ay_j\to M$，并选取诸 $\phi(y_j)\in M$ 的提升 $g_j\in\oplus Ax_i$。那么我们可以写出
\begin{align*}
M &= \oplus Ax_i \,/\, \bigl(\beta(f_1),\dots,\beta(f_r)\bigr)\\
&= \bigl(\oplus Ax_i\oplus Ay_j\bigr) \,/\, \bigl(\beta(f_1),\dots,\beta(f_r),\, y_1-g_1,\dots,y_p-g_p\bigr).
\end{align*}
由于诸 $\{\phi(y_j)\}$ 生成 $M$，对每个 $i=1,\dots,n$，我们可以选取 $h_i\in\bigoplus_{j=1}^{p}Ay_j$（也就是说，一个映射 $h\colon\bigoplus_{i=1}^{n}Ax_i\to\bigoplus_{j=1}^{p}Ay_j$），使得 $\alpha(x_i)=\phi(h_i)$。于是
\begin{align*}
M &= \bigl(\oplus Ax_i\oplus Ay_j\bigr) \,/\, \bigl(\beta(f_1),\dots,\beta(f_r),\, y_1-g_1,\dots,y_p-g_p,\\
&\qquad\qquad\qquad x_1-h_1,\dots,x_n-h_n\bigr)\\
&= \oplus Ay_j \,/\, \bigl(h(\beta(f_1)),\dots,h(\beta(f_r)),\, y_1-h(g_1),\dots,y_p-h(g_p)\bigr).
\end{align*}
\hfill$\square$

\textbf{6.4.B.（练习：有限表现满足\trAffineCommunicationLemma{}~5.3.2 的假设）}\quad

(a)~设 $f\in A$。若 $M$ 是一个\trFinitelyPresented{}的 $A$-\trModule，证明：$M_f$ 是一个\trFinitelyPresented{}的 $A_f$-\trModule。

(b)~设 $(f_1,\dots,f_n)=A$，且对一切 $i$，$M_{f_i}$ 都是\trFinitelyPresented{}的 $A_{f_i}$-\trModule。证明：$M$ 是一个\trFinitelyPresented{}的 $A$-\trModule。（提示：先用练习~6.4.A(b) 证明 $M$ 是\trFinitelyGenerated{}的，于是可写 $M=\coker(\alpha\colon N\to A^{\oplus n})$，其中 $N\subset A^{\oplus n}$ 是某个\trSubmodule。我们希望证明 $N$ 是\trFinitelyGenerated{}的。对诸 $f_i$ 作\trLocalization，用引理~6.4.6（「有限表现蕴涵总是有限表现」）证明 $N_{f_i}$ 是\trFinitelyGenerated{}的 $A_{f_i}$-\trModule，再把练习~6.4.A(b) 用到 $N$ 上。）

\textbf{6.4.C.（练习：\trCoh{}性满足\trAffineCommunicationLemma{}~5.3.2 的假设）}\quad

(a)~设 $f\in A$，而 $M$ 是一个\trCoh{}的 $A$-\trModule。证明：$M_f$ 是一个\trCoh{}的 $A_f$-\trModule。

(b)~设 $(f_1,\dots,f_n)=A$，且对一切 $i$，$M_{f_i}$ 都是\trCoh{}的 $A_{f_i}$-\trModule。证明：$M$ 是一个\trCoh{}的 $A$-\trModule。（提示：若 $\phi\colon A^{\oplus p}\to M$，则 $(\ker\phi)_{f_i}=\ker(\phi_{f_i})$，后者对每个 $i$ 都\trFinitelyGenerated。然后应用练习~6.4.A(b)。）

\textbf{6.4.7.\ \trhDefinition{}。}\quad 称一个\trQCSheaf{} $\scr{F}$ 是\textbf{\trFiniteType}的（相应地，\textbf{\trFinitelyPresented}、\textbf{\trCoh}的），若对每个\trAffine\trOpenSet{} $\Spec A$，$\Gamma(\Spec A,\scr{F})$ 都是一个\trFinitelyGenerated{}的（相应地，\trFinitelyPresented、\trCoh{}的）$A$-\trModule。注意，\trCoherentSheaf{}总是\trFiniteType{}的，而且在\trLocally\trNoetherian{}概形上，这三个概念是同一个（由命题~6.4.2）。命题~6.4.3 蕴涵：$X$ 上的诸\trCoherentSheaf{}构成一个\trAbelianCategory，我们把它记作 $\catCoh_X$。

由\trAffineCommunicationLemma{}~5.3.2 以及练习~6.4.A、6.4.B 与 6.4.C，只需在单独一个仿射\trCover{}中的诸\trOpenSet{}上检验「\trFiniteType{}性」（相应地，有限表现、凝聚性）即可。

\textbf{6.4.8.}\ \emph{\trhWarning{}。}\quad 在后来的文献中，把\trCoh{}错误地定义为\trFinitelyGenerated，这种事并不罕见。请只使用正确的定义，因为错误的定义会招致混淆。你这样做，除了出于诚实的理由，还能帮你看清：要证明那些结论，究竟哪些假设是真正必需的。而这一点总能帮你记住证明本身是什么——从而记住结论为何为真。

\textbf{6.4.9.}\ \emph{为什么要\trCoh{}？}\ 命题~6.4.3 为\trCoh{}性的定义提供了一个好动机：它给出了一个小的（此处并非专业意义上的）\trAbelianCategory，让我们可以在其中思考\trVectorBundle。

有两类人应当关心这个定义的细节，而不肯安居于「\trCoh{}即\trFiniteType」的\trNoetherian{}世界。复几何学家应当关心。他们考虑的是带\trClassicalTopology{}的\trComplexAnalyticSpace。可以用与此处类比的方式来定义\trCoh{}的 $\OO_X$-\trModule{}这一概念（见 [Se1, Def.~2]）。于是，\trOka{}（Oka）定理断言：$\mathbb{C}^n$ 的\trStructureSheaf{}（从而任何\trComplex\trManifold{}的\trStructureSheaf）是\trCoh{}的，而这是一件非常困难的事（见 [GR, \S2.5] 或 [Rem, \S7.2]）。

第二类应当关心的人，是那类可能需要同非\trNoetherian{}环打交道的做算术的人——见 \S3.6.21——或者是在非\trNonArchimedean{}解析几何中工作的人。

\textbf{6.4.10.}\ \emph{注记：\trRingedSpace{}上的\trQCSheaf{}与\trCoherentSheaf。}\quad 我们将要讨论的是概形上的\trQCSheaf{}与\trCoherentSheaf，但它们都可以在更一般的情形下定义（\trQCSheaf{}见练习~6.3.B；\trCoherentSheaf{}见 [Se1, Def.~2]）。我们所陈述的结果中，有许多在更大的一般性下依然成立，但由于证明看上去略有不同，我们把自己限制在概形上，以免分心。

\textbf{6.4.11.}\ \extrahard\ \emph{\trCoh{}在\trSmooth{}几何中不是一个好概念。}\quad 下面这个出自 B.~\trConrad{}的例子表明：在十分合理（却较不「刚性」）的情境中，\trStructureSheaf{}在自身之上并不\trCoh。考虑由 $0\in\mathbb{R}$ 处\trSmooth{}（$\mathrm{C}^{\infty}$）函数的\trGerm{}组成的环 $\OO_0$，坐标记为 $x$。$\OO_0$ 是一个\trLocalRing。它的\trMaximalIdeal{} $\mathfrak{m}$ 由 $x$ 生成。（关键想法：设 $f\in\mathfrak{m}$，并设 $f$ 有一个定义在 $(-\epsilon,\epsilon)$ 上的代表元。于是对 $t\in(-\epsilon,\epsilon)$，有 $f(t)=\int_0^t f'(u)\,du=t\int_0^1 f'(tv)\,dv$。反复「在积分号下求导」，便可验证 $\int_0^1 f'(tv)\,dv$ 是\trSmooth{}的。）

设 $\phi\in\OO_0$ 是这样一个\trSmooth{}函数的\trGerm：它在 $x\le 0$ 时为 $0$，在 $x>0$ 时为正（例如在 $x>0$ 时 $\phi(x)=e^{-1/x^{2}}$）。考虑映射 $x\phi\colon\OO_0\to\OO_0$。其\trKernel{}是\trIdeal{} $I_\phi$，由在 $x\ge 0$ 时消失的那些函数组成。显然，$I_\phi$ 非零（例如 $\phi(-x)\in I_\phi$）；但由于 $\mathfrak{m}=(x)$，有 $I_\phi=xI_\phi$，故 $I_\phi$ 不可能\trFinitelyGenerated，否则将与\trNakayama{}引理~8.2.9 矛盾。（本质上相同的论证表明：$\mathbb{R}$ 上的\trSmooth{}函数\trSheaf{}不是\trCoh{}的。）这正是\trCoh{}性在\trSmooth\trManifold{}上没有有用含义的原因。
